\section{Original partner paths at general strong-feed equality}
\label{sec:weak-equality-original-partner-paths}

Let $Q$ be a terminal weak protection closure of a finite nonempty
oriented graph $D$, using current FORCED rounds and safe single WEAK
steps. Let $T$ be a maximum-score reference completion of $Q$, and
let $L$ be a global median order of $T$ with feed $f$.
Suppose $g_D(f)>0$ and the strong feed inequality is an equality:
\[
 |G_L|=|N_T^+(f)|.
\]
All terminal remaining directions are free and forward in every median.
The original first set $O=N_D^+(f)$ and exact second set at every
closure stage are unchanged. Write
\[
 R=\mathcal M_Q(f),\quad
 R_2=R\cap N_D^{++}(f),\quad
 R_\infty=R\cap U_T(f).
\]
Equality sedimentation and free-arc rigidity exclude $R\cap G_L$.
Every newly reached original-far head belongs to $G_L$, since its
second arc is free and forward. Consequently
\[
 R=R_2\,\dot\cup\,R_\infty,
 \qquad R_2\subseteq N_D^{++}(f)\setminus G_L,
 \qquad R_\infty\subseteq\mathcal M_D(f)\cap U_D(f).
 \tag{WP1}\label{eq:weak-equality-partner-partition}
\]

\begin{theorem}[An original seed-path or unpaid-partner cycle]
\label{thm:weak-equality-original-partner-path-cycle}
Let $X$ consist of the vertices reachable from $R_2$ in the original
induced graph $D[R]$, including the seed vertices themselves, and put
$U=R_\infty\setminus X$. Then either $U=\varnothing$, or the original
graph $D[U]$ contains a directed cycle of length at least three.
Moreover, weak closure creates no new reachability from $R_2$ into
$R_\infty$ within the fixed partner set $R$:
\[
 \operatorname{Reach}_{Q[R]}(R_2)
   =\operatorname{Reach}_{D[R]}(R_2).
 \tag{WP2}\label{eq:weak-equality-original-seed-reachability}
\]
This statement needs no minimality or residual hypothesis.
\end{theorem}
\begin{proof}
Hold the final set $R$ and its partition fixed throughout the earlier
closure stages. Consider a FORCED step adding an internal arc
$a\to b$ with $b\in R_\infty$, and let
$a\triangleleft_H w$, $w\to b$ be its preceding-graph witness.
Since $b$ is far from $f$ in the final tournament, $w\to b$ forces
$w\to f$ there. If $wf$ were adjacent in $Q$, its direction would
already be $w\to f$. The retained relation $a\triangleleft w$
would then force the still-missing $a\to f$ at termination.
Thus $w\in R$.

Suppose $a$ was already reachable from the seeds in $H[R]$.
If $a\notin R_2$, a seed path has a final old arc $p\to a$ in $R$;
incoming protection nesting gives $p\to w$ in $H$, so $b$ was already
reached by the old path ending $p\to w\to b$.
If $a\in R_2$, choose an original first intermediate $o\in O$ with
$o\to a$. Nesting gives $o\to w$ in $H$. As $wf$ stays missing,
$w$ is an exact second neighbor of $f$ in this stage and hence an
original exact second neighbor. Thus $w\in R_2$ and its old arc
$w\to b$ already reaches $b$ from a seed.

For a WEAK step adding internal $a\to b$ with $b\in R_\infty$,
the old containment is $N_H^-(a)\subseteq N_H^-(b)$.
If a seed path to $a$ has positive length and last predecessor $p$,
the old arc $p\to b$ already reaches $b$. If $a\in R_2$, its original
first predecessor $o\to a$ would force $o\to b$ before the step,
contradicting that $b$ is far from $f$ in the retained final tournament.
So this case cannot occur. An added arc into $R_2$ cannot add seed
reachability because its head is already a seed. Steps outside $R$
do not alter the induced graph. Induction proves
\eqref{eq:weak-equality-original-seed-reachability}; simultaneous
FORCED additions use their common preceding graph and obey the same
argument.

For $r\in U$, the feed is maximal in terminal protection and $rf$
remains missing. Absence of the eligible WEAK direction $r\to f$
provides $p\in N_Q^-(r)\setminus N_Q^-(f)$. Since $r$ is far in $T$,
$f\to p$ in $T$ would produce a two-step path to $r$. Thus $p\to f$
in $T$, but not in $Q$, so $p\in R$. It cannot be reached from
$R_2$ in $Q[R]$, or its arc $p\to r$ would reach $r$ too.
By \eqref{eq:weak-equality-original-seed-reachability}, $p\in U$.
Therefore every vertex of $Q[U]$ has an incoming arc inside $U$,
and if $U$ is nonempty the finite oriented graph $Q[U]$ has a cycle.

To lift it to $D[U]$, consider an internal FORCED arc $a\to b$ of
$U$. The argument above supplies a witness $w\in R$.
If $w\in X$, the old arc $w\to b$ reaches $b$ from the seeds,
contrary to $b\in U$. Hence $w\in U$: the new arc shortcuts an
old internal path $a\to w\to b$ and cannot make an acyclic induced
graph cyclic. For a single WEAK internal addition, a new cycle would
have an old path $b\leadsto p\to a$. Old predecessor containment
gives $p\to b$ as well, so the old graph already had a cycle.
Steps outside $U$ do not alter its induced arcs. An acyclic $D[U]$
would therefore remain acyclic, contradicting the terminal cycle.
Loops and directed two-cycles are absent, so its length is at least
three.
\end{proof}

\subsection{Residual two and three localize the nongood obstruction}

Now assume the original $D$ is strongly connected, all-deficient and
minimal under arbitrary nonempty arc-set deletion. All packet variables
in what follows are computed in $D$. Write $\eta=\eta_D(f)$,
$g=g_D(f)$, $t=t_D(f)$ and $p=|P_D(f)|$.
For $p>0$ use the original partner packet with exceptional missing set
$B$, count $k$, deleted-outgoing count $s$, and remainder $L_D$:
\[
 |L_D|\le\eta-2k-g\mathbf1_{s>0},\qquad
 \mathcal M_D(f)\cap U_D(f)\subseteq B\cup L_D,
\]
and the original shared far-head budget for every nonexceptional
missing partner. An exceptional far partner may belong to the packet
remainder through $\partial^2P_D(f)$. No general exclusion of that
possibility is assumed: a nonexceptional missing predecessor pointing
to an exceptional far partner merely places that head in $L_D$ and
charges one unit of the actual remainder budget.

\begin{corollary}[Low residual forces an original nongood-second seed]
\label{cor:weak-equality-low-residual-seed}
Under strong feed equality, if $\eta\in\{2,3\}$, then
\[
 R_2\ne\varnothing,
 \qquad R_\infty\subseteq\operatorname{Reach}_{D[R]}(R_2).
\]
More specifically, at residual two $R_\infty=\varnothing$.
At residual three $|R_\infty|\le1$; if it consists of a vertex $r$,
some $q\in R_2$ has the original arc $q\to r$.
At residual three the sole far partner is not asserted nonexceptional.
\end{corollary}
\begin{proof}
If $p=0$, the residual identity
$\eta=2\mu_D(f)+g-1$, with $g\in\{1,2\}$ and
$\eta\in\{2,3\}$, gives $\mu_D(f)=1$. Hence $|R|\le1$.
If $p>0$, nonnegativity of the packet remainder gives $k\le1$.
When $k=0$, the original far-head packet gives $t\le\eta\le3$,
and the strong-equality missing-degree bound gives
$|R|\le t-g\le2$. When $k=1$, the original missing-far budget gives
$|R_\infty|\le\eta-k\le2$. Thus in every branch the unreachable
set $U$ of the theorem has at most two vertices; it cannot contain
the forced original cycle. Therefore $U$ is empty. Since $R$ is
nonempty, $R_2$ must be nonempty as well.

At residual two, the $p=0$ or $k=0$ bounds give $|R|\le1$,
so the nonempty seed set consumes $R$ and $R_\infty$ is empty.
For $k=1$, $L_D=\varnothing$. Every originally missing far vertex
is consequently the sole exceptional partner $b$, if one exists.
If $b\in R_\infty$, its original seed path has a final arc
$q\to b$ from another partner, necessarily nonexceptional.
The original shared head budget would place $b$ in
$L_D=\varnothing$, a contradiction.
So $R_\infty$ is empty in this branch too.

At residual three, the $p=0$ branch gives no far partner and the
$k=0$ bound $|R|\le2$ gives at most one, because there is a seed.
If $k=1$ and there were two far partners, at least one is a
nonexceptional missing far head and itself belongs to $L_D$, whose
size is at most $3-2=1$. The other must therefore be the unique
exceptional partner $b$. Its original seed path has a last predecessor
distinct from $b$, hence nonexceptional. The shared head budget places
$b$ in $L_D$ too, giving two distinct heads in a one-element budget,
a contradiction. Thus there is at most one far partner in this branch
as well; that partner may be exceptional. With only one far partner $r$, a
simple original path from the seeds to $r$ has its last predecessor
in $R_2$, giving the claimed original arc $q\to r$.
\end{proof}

\begin{corollary}[The nonexceptional residual-three far-partner shape]
\label{cor:weak-equality-eta-three-nonexception-shape}
If $\eta=3$, $p>0$, $k=0$ and $R_\infty\ne\varnothing$, then
\[
 g=1,\quad t=3,\quad s=0,\quad |R|=2,\quad
 |R_2|=|R_\infty|=1,\quad |A_f|=2,\quad g_T(f)=-1.
\]
Writing $R_2=\{q\}$ and $R_\infty=\{r\}$, the original arc
$q\to r$ exists. The full original packet also satisfies
$N_D^{--}(f)=\partial P_D(f)$ and
$L_D=\mathcal F_D(f)$ with $|L_D|=3$.
\end{corollary}
\begin{proof}
There is a seed and a far partner, so $|R|\ge2$.
The bounds $|R|\le t-g\le3-g$ force $g=1$, $t=3$, $|R|=2$,
and also $s=0$ in the stronger packet bound $t\le3-g\mathbf1_{s>0}$.
The far partner is not a new head, so $|A_f|\le t-1=2$.
The old nongood seed consumes one nongood old second head; strong
equality gives $|A_f|-g\ge1$, so $|A_f|=2$ and completed gap is
$g-|A_f|=-1$. The preceding corollary gives $q\to r$.
The partition has exceptional/outgoing defect $k+s=0$, so the incoming
packet is full. The inclusion $\mathcal F_D(f)\subseteq L_D$ and
the three-element far budget saturate $|L_D|\le3$.
\end{proof}

\begin{corollary}[The exceptional-count-one far-partner shape]
\label{cor:weak-equality-eta-three-one-exception-remainder}
If $\eta=3$, $p>0$, $k=1$ and $R_\infty\ne\varnothing$, then
$s=0$, $|L_D|=1$ and the sole terminal far partner $r$ belongs to
$L_D$. It may be exceptional or nonexceptional, and some original
nongood-second terminal partner $q$ has $q\to r$ in $D$.
\end{corollary}
\begin{proof}
The preceding classification gives the original arc $q\to r$.
If $r$ is nonexceptional, its original missing-far status puts it in
$L_D$ directly. If $r$ is exceptional, it is the unique exceptional
partner and $q\ne r$ is nonexceptional; the shared head budget again
puts $r$ in $L_D$. Hence $|L_D|\ge1$, while the packet estimate gives
$|L_D|\le1-g\mathbf1_{s>0}$. It follows that $s=0$ and
$L_D=\{r\}$.
\end{proof}

\paragraph{Scope and remaining obstruction.}
The seed vertices are original nongood \emph{second} targets, not
original far targets and not necessarily actual BAD sources. These
results do not eliminate the residual-two all-second trap or the
residual-three second-to-far path. In particular they do not assert
an original cycle on all remaining partners at negative completed gap.
They isolate the precise original path that replaces the zero-gap
unpaid-partner cycle, without applying minimality to an altered graph.

\paragraph{A corrected packet-layer assumption.}
An earlier draft incorrectly excluded exceptional far heads from
$\partial^2P_D(f)$ and hence from $E$. That exclusion is not used
in this corrected proof. The strong six-vertex geometry control
$[34,4,24,16,1,2]$ at $f=0$ has
\[
 P=\{1,5\},\ C=\{2\},\ B=\{3\},\quad
 \partial^2P=\{3,4\},\quad E=\{2,3,4\},\quad L=\{3,4\}.
\]
The exceptional missing far head $3$ therefore does belong to the
external second layer and to $L$. Other vertices have nonpositive gaps,
so this control does not meet MIN and its numerical packet budget is
not applied. It verifies the geometric scope limit only; it is not
an SNC counterexample. The corrected residual-three conclusion allows
its sole terminal far partner to be exceptional and charges it to the
one-element remainder rather than forbidding it.
